% Propietario: output/LEY9_SINTESIS_300_SUCESORA_SIN_REGRESIONES_20260904/
% manuscrito/sections/autonomia/03_prolongacion_tpk.tex:354--496.
% Microetapa CG-007. Exposicion autocontenida de dominios y prueba.
\section{Frontera coinductiva y estabilización de prolongaciones compatibles}
\label{scope:g9-frontera}

La emisión inicial tiene imagen $\mathcal W_6$ de cardinal $243$ dentro
del espacio de $729$ palabras de seis trits. Sus prolongaciones conservan
el estado enriquecido: palabra, posición, dirección, hoja, régimen,
cocientes, acarreo, frontera, incidencias y registro cronológico. Para
$r\ge1$, denotamos por $\widehat{\mathcal W}_{6r}$ el conjunto de estados
admisibles de longitud $6r$ efectivamente obtenidos mediante las
transiciones del TPK, con sus prefijos compatibles y su coordenada de
supervivencia. La proyección de palabra toma valores en
$(\mathbb F_3^6)^r$.

El truncamiento $t_{r+1,r}$ retiene el prefijo enriquecido de longitud
$6r$. La relación de prolongación es
\[
 \mathcal P_r=\{(x,y):x\in\widehat{\mathcal W}_{6r},\
 y\in\widehat{\mathcal W}_{6(r+1)},\ t_{r+1,r}(y)=x\}.
\]
Las relaciones utilizan los estados admisibles producidos por el TPK;
la igualdad de prefijos expresa su compatibilidad. Para $s>r$, escribimos
$t_{s,r}=t_{r+1,r}\circ\cdots\circ t_{s,s-1}$. La conservación de los
campos exige $t_{u,r}=t_{s,r}t_{u,s}$ cuando $u>s>r$.

\begin{proposition}[Restitución del antecedente por truncamiento]
Si $p_r$ es una sección de prolongación cuyo grafo está contenido en
$\mathcal P_r$, entonces $t_{r+1,r}p_r=I$. Para una sucesión de tales
secciones, la composición de sus truncamientos recupera el estado
antecedente completo.
\end{proposition}
\begin{proof}
La inclusión del grafo significa, para cada $x$ en el dominio de la
sección, $t_{r+1,r}(p_r(x))=x$. Al componer dos secciones, se cancela
primero la pareja de mayor profundidad y después la precedente. La
inducción concluye la fórmula para cualquier número finito de etapas.
La igualdad se refiere a los estados enriquecidos, por lo que conserva
todos sus campos y sus relaciones de prefijo.
\end{proof}

\subsection{Candidatos en la primera frontera y relación de apertura}

Sea $\widetilde{\mathcal R}_{36}$ el conjunto de estados admisibles de
profundidad $36$ con registro de frontera y memoria de borde definidos.
Para $x\in\widetilde{\mathcal R}_{36}$, el conjunto
$\operatorname{Cand}_6(x)$ consta de los estados $y$ de profundidad
$42$ obtenidos por una trayectoria tipada
\[
 x=x_{36}\xrightarrow{\mathcal U_{36}}x_{37}
 \xrightarrow{\mathcal U_{37}}\cdots
 \xrightarrow{\mathcal U_{41}}x_{42}=y,
 \qquad \tau_{42,36}(y)=x.
\]
La cola visible de cada candidato tiene seis trits y, por tanto, su
proyección de palabra tiene a lo sumo $729$ valores. La cardinalidad de
esa proyección se distingue de la cardinalidad de los estados con
registros completos.

El predicado $g_{9,x}(y)$ exige que todos los prefijos sean compatibles
y que existan continuaciones admisibles a profundidad arbitraria. La
frontera efectiva y su apertura se definen por
\begin{align*}
 \mathcal R_{36}
 &=\{x\in\widetilde{\mathcal R}_{36}:
      \exists y\in\operatorname{Cand}_6(x),\ g_{9,x}(y)=1\},\\
 G_9&=\{(x,y):x\in\mathcal R_{36},\
                y\in\operatorname{Cand}_6(x),\ g_{9,x}(y)=1\}.
\end{align*}
La construcción conserva el cuadrado de truncamiento
$\tau_{42,36}(y)=x$ y da
\[
 G_9\subseteq\mathcal U_{41}\circ\cdots\circ\mathcal U_{36}
 \subseteq\widehat{\mathcal X}_{36}\times\widehat{\mathcal X}_{42}.
\]
Cada elemento de $\mathcal R_{36}$ posee un candidato superviviente por
su definición. La relación $G_9$ abre un bloque de seis trits; la
holonomía $\Gamma_9$ compone nueve transiciones elementales. El primer
objeto es una relación entre niveles de profundidad y el segundo es el
retorno del calendario con su incremento de memoria.

\subsection{Existencia de una historia completa}

\begin{theorem}[Estabilización con conservación de la memoria]
Supóngase que los prefijos admisibles proceden de un conjunto finito
no vacío de condiciones iniciales y forman un sistema finitamente
ramificado, con estados a toda profundidad, que sus truncamientos
satisfacen la compatibilidad anterior y conservan los campos del estado.
Supóngase asimismo la naturalidad del transporte nonádico respecto del
truncamiento y las identidades de retorno $g_{k+9}=g_k$,
$m_{k+9}=m_k+1$ sobre las trayectorias admitidas. Entonces existe una
historia infinita compatible. A lo largo de ella, cada retorno conserva
la fase y aumenta la memoria entera en una unidad.
\end{theorem}
\begin{proof}
Se organiza cada prefijo por su longitud y se une a sus prolongaciones
inmediatas. Si hay varios estados iniciales, se añade una raíz formal
con ese conjunto finito de sucesores. La no vacuidad a toda profundidad
da extensiones de longitud arbitraria desde esa raíz.

Un vértice que admite extensiones arbitrariamente largas tiene al menos
un sucesor con la misma propiedad. En efecto, si cada uno de sus
finitos sucesores tuviera una cota finita de profundidad, el máximo de
esas cotas, aumentado en uno, acotaría todas las extensiones del
vértice. Esto contradice la propiedad inicial. Se escoge sucesivamente
un sucesor con extensión ilimitada; el orden finito de las palabras y
de los registros permite fijar la elección cuando se dispone de esa
codificación. La iteración produce una cadena con un prefijo en cada
nivel.

La relación de prolongación garantiza que cada truncamiento inmediato
restituye el antecedente. Su compatibilidad por composición demuestra
la misma igualdad entre cualquier par de profundidades. La naturalidad
permite aplicar la identidad de retorno a todos los prefijos donde esté
definida. Sus dos componentes proporcionan, respectivamente, la fase
conservada y el incremento unitario de memoria. Tras $s$ vueltas,
$g_{k+9s}=g_k$ y $m_{k+9s}=m_k+s$, por inducción.
\end{proof}

La prueba establece existencia a partir de ramificación finita,
prolongación a toda profundidad y compatibilidad. El predicado de
supervivencia expresa esa propiedad de la trayectoria completa. La
individualización de una sección regional concreta utiliza además el
lector y la regla funcional de esa región, desarrollados en su cadena
propia. Estos dos resultados conservan sus respectivos dominios y se
componen cuando la sección seleccionada satisface las relaciones de
prolongación anteriores.
